\clearpage
\section{プレビュー（写真の表示）}\label{sec:preview}

\begin{figure}[H]
\centering
\begin{screenshot}[0.82\linewidth]{canvas}
  \calloutabove{1}{autostretch}
  \calloutabove{2}{filename}
  \calloutabove{3}{zoomout}
  \calloutabove{4}{zoomin}
  \calloutabove{5}{zoomreset}
\end{screenshot}
\caption{プレビュー}
\end{figure}

\begin{callouts}
\co{1} \ui{Auto Stretch}：暗い写真を見やすい明るさにして表示します。\textbf{表示が変わるだけで、合成や書き出しの結果は変わりません。}
\co{2} いま表示している写真の名前です。合成した結果を表示しているときは「スタック結果」と表示されます。
\co{3} \ui{−}：表示を縮小します。
\co{4} \ui{＋}：表示を拡大します（最大 1000\%）。
\co{5} \ui{1:1}：拡大・移動をやめて、写真全体が入る大きさに戻します。
\end{callouts}

\subsection{マウスとトラックパッドの操作}

\begin{center}
\small
\begin{tabularx}{\linewidth}{l X}
\toprule
したいこと & 操作 \\
\midrule
拡大・縮小 & トラックパッドでピンチ（2本の指を広げる・つまむ）、または\key{⌘}キーを押しながらスクロール \\
拡大した写真を動かす & 写真をつかんでドラッグ \\
ブラシで塗っているときに写真を動かす & \key{space}キーまたは\key{option}キーを押しながらドラッグ \\
\bottomrule
\end{tabularx}
\end{center}
